\section{Defensa}

 Seguridad, defensa y soberanía digital en Colombia
Resultados del Ministerio de Defensa, el Ejército Nacional, la Fuerza Aeroespacial Colombiana, la Armada Nacional y la Policía Nacional en materia digital 
La soberanía digital en Colombia encuentra en el sector defensa uno de sus espacios más concretos de materialización, en la medida en que la capacidad estatal no se limita a la formulación normativa, sino que se traduce en control efectivo sobre infraestructuras, datos y dinámicas de riesgo en el ciberespacio. Desde la perspectiva de Robert Keohane y Joseph Nye, este proceso puede entenderse como un intento por gestionar la interdependencia global en términos de vulnerabilidad y capacidad de ajuste. En un entorno caracterizado por flujos constantes de información y amenazas transnacionales, el problema no es la conexión en sí misma, sino los costos a la exposición diferencial \parencite{keohane_nye_power_2011,mintic_seguridad_digital_2025}.

En este contexto, las Fuerzas Militares y la Policía han desarrollado una arquitectura de ciberdefensa basada en monitoreo, prevención y respuesta que refleja un tránsito desde esquemas reactivos hacia modelos anticipativos de gestión del riesgo. La implementación de sistemas SIEM para análisis predictivo, soluciones WAF para la protección de aplicaciones críticas, firewalls dinámicos para segmentación de tráfico y tecnologías Anti-DDoS de nueva generación respaldadas por inversiones superiores a 9.006 millones de pesos en herramientas de seguridad y más de 4.142 millones de pesos en capacidades de monitoreo y forense digital no solo amplía la capacidad técnica del Estado, sino que reduce su vulnerabilidad estructural frente a eventos externos. En términos de interdependencia compleja, estas capacidades no eliminan la exposición a amenazas globales, pero sí disminuyen los costos asociados a ellas al mejorar la capacidad de respuesta y adaptación \parencite{mindefensa_audiencia_rendicion_2025,cogfm_informe_ejecutivo_2025}.
Este proceso se profundiza con el desarrollo de tecnologías de decepción cibernética, con inversiones cercanas a 2.369 millones de pesos, orientadas a la creación de entornos simulados que permiten identificar patrones de comportamiento malicioso.  \parencite{mindefensa_audiencia_rendicion_2025}.

Sin embargo, estos avances no pueden entenderse únicamente desde la lógica de la interdependencia. Desde la perspectiva de Fernando Henrique Cardoso y Enzo Faletto, el problema central no es solo la exposición a riesgos externos, sino la posición estructural que ocupa el país en el sistema tecnológico global. En este sentido, la construcción de capacidades de ciberdefensa puede interpretarse también como un intento de reducir la dependencia tecnológica, en la medida en que fortalece el control nacional sobre infraestructuras críticas y procesos de decisión. No obstante, esta reducción es parcial: las tecnologías utilizadas, los estándares y gran parte de los ecosistemas digitales continúan anclados en redes globales dominadas por actores externos, lo que sugiere que la soberanía digital se configura como una autonomía relativa más que como independencia plena \parencite{cardoso_faletto_dependencia_1977}.

Estas capacidades técnicas se articulan con una gobernanza institucional ampliada del ciberespacio, caracterizada por la coordinación entre múltiples entidades estatales. La activación del PMU Ciber electoral, con la participación de ocho entidades, evidencia la creciente centralidad de la seguridad digital en la protección de procesos democráticos y confirma la existencia de múltiples canales de interacción uno de los rasgos centrales de la interdependencia compleja en los que actores estatales, agencias especializadas y organismos internacionales participan simultáneamente en la gestión del riesgo \parencite{mindefensa_audiencia_rendicion_2025,mintic_seguridad_digital_2025}.

De manera complementaria, el Centro Cibernético Policial ha consolidado una capacidad operativa significativa, reflejada en 17.321 bloqueos preventivos de páginas web (12.210 asociadas a abuso sexual infantil y 5.111 a juegos ilegales), la atención de 11.373 incidentes informáticos mediante el CAI Virtual, la emisión de 245 alertas preventivas en redes sociales, la identificación de 67 riesgos cibernéticos y la detección de 2.149 intentos de phishing dirigidos a entidades públicas y privadas. Estas acciones se complementan con operaciones bajo la estrategia ESCIB--C2, que han derivado en 46 capturas (31 por delitos informáticos y 15 por explotación sexual infantil). Además ha desarrollado el desarrollo del Plan Pescador, que ha permitido la incautación de dispositivos, el análisis forense digital y la judicialización de redes criminales relacionadas a la explotación infantil en entornos digitales en cooperación con agencias internacionales como Homeland Security Investigations (HSI) y el Centro Nacional para Niños Desaparecidos y Explotados (NCMEC) \parencite{policia_rendicion_logros_retos_2025,policia_plan_pescador_ninez_digital}.

Más allá de su dimensión operativa, este conjunto de acciones evidencia que la soberanía digital no se ejerce únicamente mediante infraestructura tecnológica, sino a través de la capacidad de ejercer autoridad efectiva en el ciberespacio. En términos analíticos, esto implica que el poder estatal no se define por el aislamiento de las redes globales, sino por la capacidad de intervenir en ellas, regular sus efectos y proteger a la ciudadanía en entornos digitales. Así, la soberanía digital en Colombia emerge como un proceso híbrido en el que la gestión de la interdependencia y la reducción de la dependencia se articulan para ampliar el margen de acción del Estado dentro de un sistema internacional altamente interconectado.

En paralelo, el sector defensa ha avanzado en la construcción de soberanía tecnológica mediante el desarrollo de software propio, un proceso que puede interpretarse como un intento de reducir la dependencia de proveedores externos y fortalecer el control estatal sobre información estratégica. herramientas como CAOCC PDF,  la plataforma ORIÓN para cesantías,  el sistema SIJEN para procesos disciplinarios,  y la modernización de SISPRO, con 108 millones de pesos. A estos desarrollos se suman SICOE 2.0 para analítica de datos en tiempo real,  y sistemas de gestión de casinos  \parencite{mindefensa_audiencia_rendicion_2025,cogfm_informe_ejecutivo_2025}.

Este proceso se complementa con la modernización del sistema de reclutamiento interoperable con la Registraduría Nacional, la implementación de la plataforma SGDEA para gestión documental interoperable entre fuerzas, el desarrollo de aplicaciones móviles como PadillApp en la Armada Nacional y la incorporación de chatbots institucionales en la Fuerza Aeroespacial Colombiana. En conjunto, estas iniciativas evidencian un tránsito progresivo hacia la producción interna de tecnología;  \parencite{arc_informe_rendicion_2023,fac_informe_rendicion_2023,cardoso_faletto_dependencia_1977}.

El fortalecimiento de la infraestructura tecnológica constituye otro eje central de esta estrategia y refuerza esta lógica de autonomía relativa. El sector defensa ha expandido sus redes de comunicación hasta alcanzar capacidades de 400 Gbps, y su capacidad de respaldo de datos de 40 TB a 223 TB, . Estas capacidades se acompañan de inversiones de 8.598 millones de pesos en redes integradas, 2.734 millones de pesos en redes de alta velocidad y 2.869 millones de pesos en sistemas de comunicaciones como FREQUENTIS, junto con la distribución de 1.849 computadores para fortalecer capacidades operativas y administrativas. En el ámbito naval, estos avances se complementan con la implementación de comunicaciones satelitales mediante sistemas Iridium, el despliegue de fibra óptica en zonas remotas, la modernización de sistemas de vigilancia marítima con inversiones superiores a 3.700 millones de pesos y la adopción de sistemas de videoconferencia institucional con inversiones cercanas a 1.498 millones de pesos \parencite{mindefensa_audiencia_rendicion_2025,arc_informe_rendicion_2023,cogfm_informe_ejecutivo_2025}.

Desde la perspectiva de la interdependencia compleja, estas inversiones no eliminan la exposición del Estado a redes globales, pero sí reducen su vulnerabilidad al mejorar su capacidad de ajuste frente a perturbaciones externas. En términos de Robert Keohane y Joseph Nye, la infraestructura digital no rompe la interdependencia, sino que la hace más manejable, al disminuir los costos asociados a fallas, ataques o interrupciones en los flujos tecnológicos. Así, la construcción de una base material robusta no implica aislamiento, sino resiliencia dentro de sistemas interconectados \parencite{keohane_nye_power_2011}.

El componente de innovación tecnológica refuerza esta transformación al evidenciar un tránsito hacia la producción de capacidades estratégicas propias. La operación del satélite FACSAT-2 permitió la captura de 725 imágenes, de las cuales 659 fueron procesadas para misiones de inteligencia y gestión del riesgo, mientras que proyectos como el Ala Zagi desarrollaron aeronaves autónomas de bajo costo para vigilancia. Paralelamente, el sector defensa consolidó una flota de 966 drones para inteligencia, reconocimiento y operaciones, desarrolló sistemas antidrones como NESHER V4 con patentes propias, ensambló sistemas de detección y neutralización e incorporó tecnologías GNSS RTK para la calibración de navegación aérea. Estas capacidades se articulan con la adopción de inteligencia artificial en procesos institucionales, el uso de analítica de datos en tiempo real mediante herramientas como Power BI, la implementación de simulaciones de realidad virtual para entrenamiento y el desarrollo de investigaciones en IA aplicada a educación y gestión institucional \parencite{fac_informe_rendicion_2023,mindefensa_audiencia_rendicion_2025,mindefensa_100_logros_2026}.

Más que un proceso de adopción tecnológica, este conjunto de iniciativas puede interpretarse como un intento de reposicionamiento del Estado dentro de la economía política de la tecnología. En términos de dependencia, implica avanzar aunque de manera parcial hacia la producción de conocimiento y capacidades propias; en términos de interdependencia, supone ampliar la capacidad del Estado para operar, negociar y ejercer control dentro de redes globales. En este sentido, la innovación tecnológica en el sector defensa no solo incrementa la eficiencia operativa, sino que redefine el lugar del Estado colombiano en el sistema internacional, al transformar su rol de usuario de tecnología en un actor con capacidades crecientes de producción y adaptación estratégica.
A estos avances se suma un fortalecimiento estructural de las capacidades estratégicas del Estado que amplía significativamente el alcance de la soberanía digital en Colombia. La Estrategia Nacional de Seguridad Digital 2025--2027 evidencia la magnitud del desafío al señalar que el país concentró aproximadamente el 17\% de los intentos de ciberataque en América Latina cerca de 36.000 millones de eventos en 2024, lo que ha llevado a la formulación de una respuesta institucional orientada a la creación de una entidad nacional especializada, la implementación de un Observatorio de Seguridad Digital y la adopción obligatoria de arquitecturas Confianza Cero (Zero Trust) en sistemas gubernamentales, junto con el fortalecimiento de los CSIRT bajo el modelo internacional SIM3. Desde la perspectiva de la interdependencia compleja, este conjunto de medidas refleja un intento por reducir la vulnerabilidad del Estado frente a flujos y amenazas transnacionales, no mediante el aislamiento, sino a través del fortalecimiento de su capacidad de anticipación, monitoreo y respuesta en redes altamente interconectadas. En este sentido, la estrategia no busca eliminar la interdependencia, sino hacerla más manejable y menos costosa \parencite{mintic_seguridad_digital_2025,keohane_nye_power_2011}.

En el plano operativo, el fortalecimiento de capacidades tecnológicas ha adquirido una dimensión más robusta mediante inversiones significativas en sistemas de defensa avanzada. Se destacan los 30.000 millones de pesos gestionados para el desarrollo de tecnologías anti-drones y la priorización de 126.200 millones de pesos en capacidades defensivas frente a amenazas aéreas no tripuladas, junto con la consolidación de una flota de 966 drones, incluyendo 66 modelos de alta capacidad que incrementaron el poder de fuego institucional en un 147\% \parencite{mindefensa_audiencia_rendicion_2025,mindefensa_100_logros_2026}. 
Paralelamente, el desarrollo del sistema NESHER V4 permitió la generación de capacidades propias de detección e inhibición, con 34 sistemas ensamblados y 35 detectores desplegados en zonas estratégicas, además de la neutralización efectiva de múltiples amenazas en campo. Estos avances reflejan una transición hacia un modelo de defensa tecnológica que integra capacidades digitales, físicas y operacionales, y que, desde el enfoque de la dependencia, puede interpretarse como un intento por reducir la subordinación tecnológica al desarrollar soluciones propias en ámbitos tradicionalmente dominados por proveedores externos. No obstante, esta reducción es parcial, en tanto dichas capacidades continúan insertas en cadenas globales de producción tecnológica \parencite{mindefensa_audiencia_rendicion_2025,cardoso_faletto_dependencia_1977}.

De manera complementaria, el Estado ha fortalecido su capacidad de gestión territorial y toma de decisiones mediante el uso de sistemas de información avanzados. Plataformas como el Sistema Integral de Derechos Humanos (SIDEH), que permitió el seguimiento de 1.399 activaciones de rutas de protección, el Sistema de Información de Gestión Territorial (SEGET), orientado a la priorización de inversiones, y el SIIPO para el seguimiento de compromisos del posconflicto, evidencian la incorporación de herramientas digitales en la gestión estatal y la articulación entre seguridad, desarrollo territorial y política pública. Desde la lógica de la interdependencia compleja, estos sistemas amplían la capacidad del Estado para procesar información y coordinar acciones en múltiples niveles, reduciendo la sensibilidad frente a perturbaciones externas al mejorar la capacidad de respuesta interna \parencite{mindefensa_siipo_implementacion_2025,mindefensa_audiencia_rendicion_2025}.

En el ámbito de ciencia, tecnología e innovación, la adopción de la política sectorial de CTeI mediante la Resolución 2443 de 2024 marca un punto de inflexión al establecer un enfoque orientado por misiones para cerrar brechas tecnológicas en la Fuerza Pública. Este proceso se refleja en la gestión de 94 procesos de propiedad intelectual y la obtención de 47 registros, así como en la implementación de esquemas de transferencia tecnológica que incluyen capacidades de mantenimiento de aeronaves, modernización de sistemas de comunicaciones y desarrollo de entrenamiento táctico avanzado. A ello se suma la exploración de tecnologías emergentes como la computación cuántica y la criptografía post-cuántica, lo que evidencia una estrategia orientada no solo a la adopción, sino a la producción de conocimiento tecnológico de frontera. En términos de dependencia, este giro resulta particularmente significativo, en tanto sugiere un desplazamiento aunque aún incipiente hacia la construcción de capacidades endógenas de innovación, elemento clave para modificar la posición estructural del país en la economía política global de la tecnología \parencite{mindefensa_audiencia_rendicion_2025,cogfm_informe_ejecutivo_2025}.

En conjunto, estos desarrollos muestran que la soberanía digital en Colombia se está configurando a través de una doble dinámica: por un lado, la gestión de la interdependencia mediante el fortalecimiento de capacidades de resiliencia frente a amenazas globales; por otro, la reducción parcial y progresiva de la dependencia mediante la construcción de capacidades tecnológicas propias. En este equilibrio, la soberanía no se expresa como independencia absoluta, sino como una forma de autonomía relativa que amplía el margen de acción del Estado dentro de estructuras globales interdependientes.

Finalmente, el fortalecimiento de capacidades en el dominio marítimo y científico amplía el alcance de la soberanía digital más allá del ciberespacio, evidenciando su carácter multidimensional. La incorporación del buque de investigación ARC Simón Bolívar, el desarrollo de sistemas de simulación inmersiva para entrenamiento, la implementación de sistemas de gestión de combate de fabricación nacional y el uso de tecnologías robóticas para investigación marina reflejan una integración progresiva entre defensa, ciencia y tecnología. Estas capacidades se complementan con el desarrollo de modelos predictivos para operaciones de búsqueda y rescate y el uso de herramientas digitales para el monitoreo ambiental y marítimo, configurando una forma de soberanía que no se limita al control de datos, sino que se extiende a la capacidad de producir, procesar y aplicar conocimiento en múltiples dominios estratégicos \parencite{arc_informe_rendicion_2023,mindefensa_audiencia_rendicion_2025}.

Desde la perspectiva de la interdependencia compleja, este tipo de desarrollos amplía la capacidad del Estado para operar dentro de redes globales altamente interconectadas, al reducir su vulnerabilidad frente a contingencias externas mediante el fortalecimiento de sus capacidades de respuesta y adaptación. Sin embargo, lejos de implicar un proceso de desconexión, estas capacidades confirman que la soberanía digital se construye en interacción con dichas redes, en un entorno donde la gestión del riesgo depende tanto de la articulación internacional como de la fortaleza interna \parencite{keohane_nye_power_2011}.

La transformación digital también se refleja en la relación entre el Estado y la ciudadanía, evidenciando que la soberanía digital no es exclusivamente una cuestión de seguridad, sino también de gobernanza. La digitalización permitió que el 98,5\% de las peticiones en la Fuerza Aeroespacial fueran gestionadas por canales virtuales, mientras que el portal institucional registró más de 575.000 visitas y las plataformas educativas alcanzaron capacidades superiores a 300.000 usuarios. A ello se suma la implementación de SALUD.SIS, con inversiones cercanas a 9.000 millones de pesos, que ha modernizado la prestación de servicios mediante telemedicina y telemonitoreo, así como la integración de bibliotecas digitales que conectan más de 60 unidades institucionales. En términos de transparencia, el sector alcanzó calificaciones de 93 sobre 100 en índices de acceso a la información, fortaleciendo mecanismos de participación digital y rendición de cuentas. Este conjunto de avances sugiere que la soberanía digital también se expresa en la capacidad del Estado para organizar, distribuir y garantizar el acceso a bienes públicos digitales \parencite{fac_informe_rendicion_2023,mindefensa_audiencia_rendicion_2025}.

En conjunto, estos desarrollos evidencian que la soberanía digital en Colombia no puede entenderse como independencia tecnológica absoluta, sino como la construcción de capacidades estatales orientadas a gestionar la interdependencia global y a reducir, de manera parcial, la dependencia estructural. Desde la perspectiva de Robert Keohane y Joseph Nye, estas capacidades disminuyen la vulnerabilidad del Estado frente a amenazas externas al mejorar su capacidad de ajuste, sin eliminar su inserción en redes globales. Desde el enfoque de Fernando Henrique Cardoso y Enzo Faletto, el fortalecimiento de software propio, infraestructura y capacidades tecnológicas representa un intento por desplazar, aunque de forma incompleta, los centros de decisión hacia el ámbito nacional, reconfigurando las relaciones de subordinación tecnológica \parencite{keohane_nye_power_2011,cardoso_faletto_dependencia_1977}.
En este equilibrio, la soberanía digital emerge como una forma de autonomía relativa, en la que el Estado amplía su capacidad de decisión, negociación y control dentro de estructuras interdependientes. Más que un proyecto de autosuficiencia, se trata de una estrategia de inserción activa en el sistema internacional, donde el poder no reside en el aislamiento, sino en la capacidad de intervenir, regular y adaptarse a las dinámicas de un entorno digital global.

% \section{Bibliografía}
% \printbibliography[heading=none]
